% KSC single-file draft. Keep kcc.cls from the Overleaf template.
\documentclass[preprint]{kcc}

\usepackage{graphicx}
\usepackage{booktabs}
\usepackage{amsmath}
\usepackage{multirow}
\usepackage{balance}

\title{출력 신뢰도와 저비용 정답 검증을 이용한 적응형 SLM--LLM 라우팅}
\author{
제1저자$^{\circ}$, 공동저자, 교신저자\\
소속기관명\\
\{first, second, corresponding\}@example.ac.kr
}
\engtitle{Adaptive SLM--LLM Routing with Output Confidence and Low-Cost Answer Verification}
\engauthor{
First Author, Second Author, Corresponding Author\\
Affiliation\\
}

\abstract{
대규모 언어 모델은 높은 추론 성능을 제공하지만 모든 요청을 대형 모델에 할당하면
계산 비용과 지연시간이 증가한다. 본 연구는 소형 모델의 질문 임베딩뿐 아니라 실제
생성 출력의 신뢰도(confidence)와 짧은 정답 전용(answer-only) 검증 결과를 이용하는 적응형 SLM--LLM 라우팅
방법을 제안한다. 제안 방법은 Qwen2.5-1.5B가 먼저 답을 생성하고, 경계 사례에 대해서만
설명을 제외한 정답 전용 검증을 수행한 후 두 숫자 답의 일치 여부에 따라
Qwen2.5-7B로 전달한다. GSM8K에서 학습한 신뢰도 라우터를 SVAMP로 전이하여
분리된 정책 선택 및 인증 분할과 공식 테스트에서 평가하였다. RTX 3090 동일 조건
실험에서 정답 전용 검증은 기존 두 번째 전체 풀이보다 문항당 지연시간을
492.97 ms에서 29.15 ms로 94.09\% 감소시켰다. 실측 지연시간으로 환산한 계단식
비용은 독립 인증 분할과 공식 테스트에서 항상 상위 모델(Always-Upper) 대비 각각 22.29\%와 19.42\%
감소했으며, 정확도는 각각 88.00\%와 85.67\%를 기록하였다. C3 기준선과의 exact
McNemar 검정에서는 유의한 정확도 차이가 관측되지 않았다. 한편 95\% 위험
상한은 5.195\%로 사전 설정한 5\% 기준을 근소하게 충족하지 못했다. 결과적으로
제안 방법은 관측된 성능--비용 frontier를 개선하지만, 안전 보장과 정확도 비열등성의
확증에는 더 큰 독립 표본이 필요하다.
}

\begin{document}
\maketitle

\section{서론}

대형 언어 모델(LLM)은 다양한 추론 과제에서 높은 성능을 보이지만, 모든 요청을 대형
모델로 처리하는 방식은 불필요한 계산 비용과 응답 지연을 유발한다. 실제 요청 중 일부는
소형 언어 모델(SLM)로도 정확히 처리할 수 있으므로, 입력별로 적절한 모델을 선택하는
라우팅은 품질과 비용을 동시에 관리하기 위한 핵심 기술이다.

기존 LLM 라우팅은 질문 임베딩, 선호 데이터 또는 별도의 분류기를 이용해 강한 모델의
필요성을 예측한다~\cite{chen2024frugalgpt,ong2025routellm}. 그러나 의미적으로 유사한
질문도 SLM 관점의 난이도가 다를 수 있으며, 질의 전용(query-only) 표현은 SLM이 실제로 생성한
답의 불확실성과 자기 일관성을 직접 반영하지 못한다. 또한 경계 사례마다 전체 풀이를
두 번 생성하면 검증 자체의 디코딩 비용이 커져 계단식 구조의 경제성이 약해진다.

본 연구는 첫 번째 SLM 출력에서 계산한 신뢰도와 별도 프롬프트 기반 정답 전용 검증을 결합한다.
라우터가 확실한 사례는 첫 답을 즉시 채택하고, 경계 사례는 동일 SLM이 숫자 답만 짧게
재생성하도록 한다. 두 답이 일치하면 SLM 출력을 채택하고, 불일치하거나 신뢰도가
낮으면 상위 LLM으로 전달한다.

본 연구의 기여는 다음과 같다. 첫째, 질문 의미와 실제 SLM 출력 신뢰도를 결합한
두 임계값 라우팅 구조를 제시한다. 둘째, 전체 재풀이 대신 지연시간 벤치마크 128문항에서
평균 5.40토큰의 정답 전용
검증을 사용해 동일 GPU에서 두 번째 호출 지연시간을 94.09\% 줄인다. 셋째, 정책 선택과
평가 분할을 분리하고 대응 부트스트랩 및 McNemar 검정을 통해 정확도--비용 상충 관계를
정량적으로 분석한다. 넷째, 안전 기준의 근소 실패도 함께 보고하여 방법의 적용 범위를
명확히 한다.

\section{관련 연구}

FrugalGPT는 이질적인 비용 구조를 가진 모델을 cascade로 결합하여 비용--품질
trade-off를 최적화하였다~\cite{chen2024frugalgpt}. RouteLLM은 인간 선호 데이터와
증강을 이용해 강한 모델과 약한 모델 사이의 라우터를 학습하였다~\cite{ong2025routellm}.
이들과 달리 본 연구는 입력 표현만으로 모델을 선택하지 않고, 약한 모델이 생성한
출력 신뢰도와 별도 프롬프트의 짧은 숫자 검증을 라우팅 신호로 사용한다.

GSM8K 연구는 생성 후보를 평가하는 verifier의 효과를 보였다~\cite{cobbe2021gsm8k}.
선택적 생성 연구는 신뢰할 수 없는 출력을 선택하거나 거부할 때 위험 통제가 필요함을
보였다~\cite{lee2024selective}. 본 연구는 별도의 검증 모델을 학습하지 않고 동일
SLM의 정답 전용 재생성을 이용해 첫 출력과의 숫자 일치를 확인한다. 전통적
self-consistency가 동일 프롬프트에서 확률적으로 여러 출력을 표본화해 다수결하는 것과
달리~\cite{wang2023selfconsistency}, 본 방법은 greedy decoding을 유지하면서 풀이
프롬프트와 정답 전용 프롬프트라는 서로 다른 두 관점을 사용하고 신뢰도 경계 사례에만
두 번째 호출을 수행한다. 따라서
별도 가중치의 적재·상주 비용 없이 지연시간 벤치마크 128문항에서 평균 5.40토큰의
추가 디코딩만 사용한다. 범용 임베딩은
의미적 유사성을 효과적으로 표현하지만~\cite{xiao2023bge}, 같은 주제의 질문도 SLM이
해결할 수 있는 난이도는 다를 수 있다. 초기 GSM8K 실험에서 frozen BGE 기반 logistic
라우터의 AUC는 0.589였고, 출력 신뢰도 기반 라우터는 0.722를 기록하였다.

\section{제안 방법}

그림~\ref{fig:architecture}는 제안 라우팅 구조를 보여준다. Lower SLM의 첫 출력으로
라우팅 점수를 계산하고, 높은 신뢰도의 요청은 즉시 채택한다. 경계 사례는 숫자
한 줄만 생성하는 별도 프롬프트 검증을 수행하며, 첫 답과 검증 답이 다르면 상위 LLM으로
전달한다.

\begin{figure}[t]
  \centering
  \includegraphics[width=\columnwidth]{figures/fig1_architecture.pdf}
  \caption{출력 신뢰도와 정답 전용 검증을 결합한 제안 라우팅 구조.}
  \label{fig:architecture}
\end{figure}

첫 출력의 11차원 특징은 token log-probability의 평균, 최솟값, 표준편차, entropy의
평균과 최댓값, top-1/top-2 probability margin의 평균과 최솟값, completion token 및
문자 수, 숫자 개수, Final answer 형식 존재 여부로 구성된다. Logistic regression은
다음 routing score를 출력한다.

\begin{equation}
q_i = \sigma(\mathbf{w}^{\top}\mathbf{x}_i+b).
\end{equation}

정책 선택 분할에서 고정한 두 임계값은 $\tau_l=0.12$, $\tau_h=0.80$이다.

\begin{equation}
\pi(q_i)=
\begin{cases}
\text{accept Lower}, & q_i \ge \tau_h,\\
\text{정답 전용 검증}, & \tau_l \le q_i < \tau_h,\\
\text{상위 모델 전달}, & q_i < \tau_l.
\end{cases}
\end{equation}

중간 영역에서는 첫 숫자 답 $a_i^{(1)}$과 검증 숫자 답 $a_i^{(v)}$가 모두 추출
가능하고 동일할 때만 Lower를 채택한다.

\begin{equation}
\text{accept}_i = [q_i\ge\tau_h]\lor
([\tau_l\le q_i<\tau_h]\land[a_i^{(1)}=a_i^{(v)}]).
\end{equation}

문항별 비용과 정규화 비용은 다음과 같다.

\begin{equation}
C_i=C_L+I_i^{(v)}C_V+I_i^{(U)}C_U, \qquad
\widetilde{C}=\frac{1}{N}\sum_{i=1}^{N}\frac{C_i}{C_U}.
\end{equation}

기준 비용은 $C_L=1$, $C_U=4.667$이며, 동일 GPU에서 측정한 정답 전용/전체 풀이
지연시간 비율에 따라 $C_V=0.0591$을 사용한다. 추가로 $C_V=0.35$인 보수적 비용
민감도도 확인한다.

\section{실험 설정}

라우터 학습에는 GSM8K를 사용하고, 교차 과제 평가는 SVAMP에서 수행하였다
~\cite{cobbe2021gsm8k,patel2021svamp}. GSM8K pilot 2,000문항은 router train 1,000,
distillation 400, validation 200, test 400으로 분리하였다. SVAMP train 700문항은
초기 검증 200, 정책 선택 250, 독립 위험 인증 250으로 겹치지 않게
분리했으며, 공식 테스트 300문항은 최종 진단에만 사용하였다. 모든 분할은 고정 seed
42와 문항 ID로 관리하였다.

Lower와 Upper는 각각 Qwen2.5-1.5B-Instruct와 Qwen2.5-7B-Instruct를 사용하였다
~\cite{yang2024qwen25}. 두 모델은 4-bit로 로드하고 단일 NVIDIA RTX 3090 24GB에서
greedy decoding으로 평가하였다. 첫 답은 최대 256토큰, 정답 전용 검증은 최대
64토큰으로 제한하였다.

평가 지표는 최종 정확도, 상위 모델 호출률, 위험 라우팅, Always-Upper 대비
정규화 비용이다. 지연시간은 워밍업 뒤 동일 128문항을 batch size 16으로 조건별
2회 ABBA 순서로 측정하였다. 통계 분석은 동일 문항에 대한 10,000회 대응 부트스트랩
신뢰구간과 정확 McNemar 검정을 사용하였다. 위험률에는 독립 250문항의 단측
95\% Clopper--Pearson 상한을 사용하였다~\cite{clopper1934use}.

비교 기준인 C3는 출력 신뢰도와 두 번째 전체 풀이의 정답 일치 여부를 결합한
선택적 일관성 정책이며, 제안 방법과 동일한 첫 번째 SLM 및 상위 LLM을 사용한다.

\section{실험 결과}

그림~\ref{fig:latency}에서 정답 전용 검증은 두 번째 전체 풀이의 평균 124.55토큰을
5.40토큰으로 줄였다. 이에 따라 문항당 지연시간은 492.97 ms에서 29.15 ms로
94.09\% 감소했고 처리량은 2.03 items/s에서 34.32 items/s로 증가하였다.

\begin{figure*}[t]
  \centering
  \includegraphics[width=0.88\textwidth]{figures/fig2_latency_tokens.pdf}
  \caption{동일 RTX 3090에서 두 번째 전체 풀이와 정답 전용 검증의 문항당
  지연시간 및 생성 토큰 수 비교.}
  \label{fig:latency}
\end{figure*}

\begin{table}[t]
\centering
\caption{SVAMP에서 고정 정책의 정확도 및 실측 지연시간 환산 비용}
\label{tab:main}
\small
\setlength{\tabcolsep}{3.2pt}
\begin{tabular}{llrrr}
\toprule
분할 & 방법 & 정확도(\%) & 위험률(\%) & 비용 \\
\midrule
인증 & Always Upper & 88.80 & -- & 1.000 \\
      & C3           & 88.80 & 2.00 & 1.008 \\
      & Answer-only  & 88.00 & 2.80 & \textbf{0.777} \\
\midrule
테스트 & Always Upper & 88.00 & -- & 1.000 \\
      & C3           & 87.67 & 1.00 & 1.045 \\
      & Answer-only  & 85.67 & 3.00 & \textbf{0.806} \\
\bottomrule
\end{tabular}
\end{table}

표~\ref{tab:main}에서 제안 방법은 Always-Upper 대비 독립 인증과 공식 테스트의
정규화 비용을 각각 22.29\%와 19.42\% 줄였다. C3 대비 정확도는 각각 0.8\%p와
2.0\%p 낮았으나 exact McNemar 검정은 유의하지 않았다. 그러나 이는 정확도 동등성
또는 비열등성을 입증하지 않는다.

\begin{figure}[t]
  \centering
  \includegraphics[width=0.92\columnwidth]{figures/fig3_pareto_frontier.pdf}
  \caption{정책 선택 분할의 성능--비용 Pareto frontier. 보라색 점은 독립 평가 전에
  고정한 임계값 0.12/0.80 정책이다.}
  \label{fig:pareto}
\end{figure}

그림~\ref{fig:pareto}는 임계값 변화에 따라 비용과 정확도가 함께 증가하는 양상을
보인다. 본 연구는 정책 선택 분할에서 품질과 위험 제약을 만족한 0.12/0.80 정책을
고정하였다.

\begin{table}[t]
\centering
\caption{정답 전용 검증과 C3의 대응 비교(단위: \%p)}
\label{tab:paired}
\small
\setlength{\tabcolsep}{2.5pt}
\resizebox{\columnwidth}{!}{%
\begin{tabular}{lrrr}
\toprule
분할 & $\Delta$정확도 [95\% CI] & $p_{McN}$ & $\Delta$비용 [95\% CI] \\
\midrule
인증 & $-0.8$ [$-2.4,0.8$] & .625 & $-23.06$ [$-28.70,-17.38$] \\
테스트 & $-2.0$ [$-4.33,0.01$] & .146 & $-23.87$ [$-29.58,-18.27$] \\
\bottomrule
\end{tabular}
}
\end{table}

독립 인증 250문항에서 정답 전용 정책의 위험 오류는 7건(2.80\%)이었다.
단측 95\% Clopper--Pearson 상한은 5.195\%로 사전 설정한 5\% 기준을 0.195\%p
초과하였다. 따라서 본 결과를 5\% 위험 보장으로 해석하지 않는다. 안전 인증을 통과한
C3 정책은 정규화 비용이 1.008로 Always-Upper보다 높았다.

본 연구에는 네 가지 제한이 있다. 첫째, 산술 문장제와 하나의 Qwen 모델 계열만
평가하였다. 둘째, 공식 테스트 300문항은 정확도 비열등성을 검정하기에 충분하지 않다.
셋째, 지연시간 측정은 단일 RTX 3090과 조건별 2회 반복에 한정되며 API 가격이나 에너지
소비와 동일하지 않다. 넷째, 정답 전용 정책은 5\% 위험 인증을 근소하게 통과하지
못했다. 후속 연구에서는 정책 선택 분할에서
$\min_{\tau_l,\tau_h}\widetilde{C}(\tau_l,\tau_h)$를 최적화하되 위험률의 단측
신뢰상한이 목표 $\alpha$ 이하라는 제약을 동시에 적용하고, 고정된 임계값을 새로운 독립
분할에서 다시 인증할 예정이다.

\section{결론}

본 연구는 SLM의 출력 신뢰도와 짧은 정답 전용 검증을 결합한 적응형 SLM--LLM
라우팅 방법을 제안하였다. 전체 풀이를 다시 생성하는 C3 방식과 비교해 정답 전용
검증은 동일 GPU에서 두 번째 호출 지연시간을 94.09\% 줄였으며, 실측 지연시간 환산
계단식 비용을 Always-Upper 대비 19.42--22.29\% 절감하였다. 동일 문항 대응
검정에서 C3와 유의한 정확도 차이는 관측되지 않았지만 이는 동등성 증명을 의미하지
않는다. 또한 95\% 위험 상한은 5.195\%로 사전 안전 기준을 근소하게 초과했다.
향후에는 더 큰 독립 표본과 다양한 모델 조합에서 비열등성, 위험 통제 및 서빙
환경의 실제 비용을 함께 검증할 예정이다.

\balance
\begin{thebibliography}{9}

\bibitem{chen2024frugalgpt}
L. Chen, M. Zaharia, and J. Zou,
``FrugalGPT: How to Use Large Language Models While Reducing Cost and Improving Performance,''
\textit{Transactions on Machine Learning Research}, 2024.

\bibitem{ong2025routellm}
I. Ong et al.,
``RouteLLM: Learning to Route LLMs with Preference Data,''
\textit{International Conference on Learning Representations}, 2025.

\bibitem{cobbe2021gsm8k}
K. Cobbe et al.,
``Training Verifiers to Solve Math Word Problems,''
arXiv:2110.14168, 2021.

\bibitem{patel2021svamp}
A. Patel, S. Bhattamishra, and N. Goyal,
``Are NLP Models Really Able to Solve Simple Math Word Problems?''
\textit{Proceedings of NAACL-HLT}, 2021.

\bibitem{yang2024qwen25}
A. Yang et al.,
``Qwen2.5 Technical Report,''
arXiv:2412.15115, 2024.

\bibitem{xiao2023bge}
S. Xiao et al.,
``C-Pack: Packed Resources for General Chinese Embeddings,''
arXiv:2309.07597, 2023.

\bibitem{lee2024selective}
M. Lee, K. Kim, T. Kim, and S. Park,
``Selective Generation for Controllable Language Models,''
\textit{Advances in Neural Information Processing Systems}, 2024.

\bibitem{clopper1934use}
C. J. Clopper and E. S. Pearson,
``The Use of Confidence or Fiducial Limits Illustrated in the Case of the Binomial,''
\textit{Biometrika}, vol. 26, no. 4, pp. 404--413, 1934.

\bibitem{wang2023selfconsistency}
X. Wang et al.,
``Self-Consistency Improves Chain of Thought Reasoning in Language Models,''
\textit{International Conference on Learning Representations}, 2023.

\end{thebibliography}
\end{document}
